\chapter*{第一部分收束：从会运行到可复现项目包}
\addcontentsline{toc}{chapter}{第一部分收束：从会运行到可复现项目包}

完成 Part I 后，你已经接触了 Linux 命令行、Git、脚本与 notebook、文件 I/O 和科学绘图。它们表面上像五组工具，但在科研训练中其实服务同一个目标：让一次分析不只是在自己电脑上“跑过一次”，而是能被自己、同伴、助教或未来的你重新找到、重新运行、重新检查。

因此，Part I 的真正交付物不是一串命令清单，而是一个\textbf{最小可复现项目包}。

\section*{五层项目包}

一个入门阶段就足够好的科研计算项目，至少应有五层结构。

第一层是\textbf{环境层}。它回答“这段代码在哪里、用什么环境运行”：

\begin{itemize}
    \item 项目目录结构清楚；
    \item 相对路径稳定；
    \item Python 版本、关键库和运行入口有记录；
    \item 临时文件、输出文件和原始数据不会混在一起。
\end{itemize}

第二层是\textbf{版本层}。它回答“这个结果对应哪一次代码状态”：

\begin{itemize}
    \item Git 提交粒度能解释一次变化；
    \item 提交前检查过 \texttt{diff}；
    \item 分支用于试错，tag 或 release 用于冻结结果；
    \item 大数据不盲目进入 Git，但数据来源、版本和校验线索被记录下来。
\end{itemize}

第三层是\textbf{运行层}。它回答“别人怎样重新运行这件事”：

\begin{itemize}
    \item notebook 用于探索、解释和展示；
    \item script 用于稳定重复执行；
    \item module 用于保存可复用逻辑；
    \item 入口脚本、参数、日志和错误信息足够清楚。
\end{itemize}

第四层是\textbf{数据层}。它回答“输入数据是否被正确理解”：

\begin{itemize}
    \item 文件路径、字段、单位、缺失值和非法值被检查；
    \item CSV 字段契约和 FITS/header 契约被当成同一类元数据约束；
    \item 数据来源、处理历史和后续分析入口有 provenance 记录；
    \item 读入成功不等于可以分析，必须先确认数据契约。
\end{itemize}

第五层是\textbf{证据层}。它回答“结果怎样变成可以讨论的科学证据”：

\begin{itemize}
    \item 图表有清楚变量、单位、坐标和图例；
    \item 对数坐标、误差棒、残差图和视觉编码都有明确用途；
    \item 图像文件本身作为科研记录保存；
    \item 图表解读遵循 claim--evidence--limit，而不是只描述形状。
\end{itemize}

\section*{Part I 的共同算法}

可以把 Part I 学到的工作方式压缩成下面这条算法：

\begin{enumerate}
    \item 建立清楚的项目目录；
    \item 记录当前环境和运行入口；
    \item 读取数据前先写字段、单位和路径契约；
    \item 用 notebook 探索，用脚本固定可重复流程；
    \item 每次重要变化前后检查 \texttt{git status} 和 \texttt{git diff}；
    \item 把结果图、日志、参数和数据版本一起保存；
    \item 用一段短说明把结果、代码版本和限制连起来；
    \item 对准备继续使用的状态打 tag 或形成 release note。
\end{enumerate}

这条算法没有复杂数学，却是后面所有数据分析、机器学习和 后续项目的地基。真正的困难不是记住命令，而是在项目变复杂、文件变多、时间拉长时仍然保持这些习惯。

\section*{一个最小交付清单}

\begin{center}
\begin{tabular}{p{0.22\textwidth} p{0.32\textwidth} p{0.32\textwidth}}
\toprule
交付对象 & 最小内容 & 最容易失效的地方 \\
\midrule
项目目录 & \texttt{data/}、\texttt{notebooks/}、\texttt{scripts/}、\texttt{results/} & 临时文件和最终结果混在一起 \\
版本记录 & 清楚提交、关键 tag、简短 release note & 结果图找不到对应 commit \\
运行入口 & notebook 解释流程，script 固定流程 & notebook 顺序依赖导致无法重跑 \\
数据记录 & 字段契约、单位、来源、缺失值处理 & 文件能读但物理含义不清 \\
图表证据 & 带单位图表、误差或残差诊断、限制说明 & 图好看但不能支撑论断 \\
\bottomrule
\end{tabular}
\end{center}

这张表可以当作每个后续小项目的提交前检查表。如果某一行完全说不清，通常说明项目还没有准备好进入下一阶段。

\section*{Part I 结束时的最小验收标准}

如果把 Part I 当作一门短课，它的结课验收不应是“背出多少命令”，而应是学生能否独立交出一个小型项目包。最小验收标准可以写成：

\begin{enumerate}
    \item \textbf{能定位}：能解释当前目录、输入数据、输出目录和脚本入口；
    \item \textbf{能运行}：能用一条命令或一个 README 说明重新运行流程；
    \item \textbf{能记录}：能保存标准输出、错误日志、关键参数和运行状态；
    \item \textbf{能追溯}：能把结果文件连接到 commit hash、数据来源和环境说明；
    \item \textbf{能检查}：能说明字段、单位、header、缺失值和质量标记如何处理；
    \item \textbf{能解释}：能用图表和 claim--evidence--limit 说明一个有限结论。
\end{enumerate}

这六条如果全部做到，学生就已经具备进入 Part II 的地基。后面即使数据类型和模型变复杂，基本判断仍然相同：结果是否能被找到、复跑、追溯、检查和解释。

\section*{从看懂到练会}

如果只阅读正文，你应该已经能理解为什么这些工具属于同一条科研工作流；如果完成配套训练，你还应该能够把这条工作流落到一个真实目录中。两者的分工不同，但最终会合在同一个交付物里。

正文要求你能解释：

\begin{itemize}
    \item 为什么路径、版本、环境、数据契约和图表解释会影响科研可信度；
    \item 为什么 notebook 适合探索，脚本适合重复运行，模块适合复用逻辑；
    \item 为什么“能打开文件”“能跑出图”“能得到模型分数”都还不等于完成分析；
    \item 为什么每个结果都应该能回到输入数据、代码版本、运行入口和解释限制。
\end{itemize}

训练材料要求你能做出：

\begin{itemize}
    \item 一个结构清楚的小型项目目录；
    \item 一组可重跑的 notebook 或脚本入口；
    \item 一份数据入口检查记录；
    \item 一张带图注、解释边界和生成记录的科学图表；
    \item 一次能说明变化意图的 Git 提交。
\end{itemize}

这就是本书后续所有章节的基本节奏：正文先建立概念地图，训练再把概念压实成操作、证据和交付。读者不需要在 Part I 结束时变成命令行专家，但应该开始具备一种科研计算直觉：任何结果都要能被找到、复跑、检查和解释。

\section*{它们最终会进入哪些证据文件}

等到后续课程项目或科研训练时，Part I 的这些基础工作不会消失，而是会进入同一套证据文件：

\begin{itemize}
    \item 项目目录、相对路径、运行入口和环境说明进入 Evidence Record；
    \item Git 提交、tag、release note 和结果版本线索也进入 Evidence Record；
    \item 数据文件位置、字段契约、来源和校验记录进入 Data Card；
    \item 图表文件、绘图代码、坐标单位和 claim--evidence--limit 说明进入 Evidence Record 或 Trust Statement；
    \item AI 助手参与脚本整理、错误解释或 README 起草时，使用记录和人工核查进入 AI Usage Log。
\end{itemize}

这也是为什么 Part I 不能只被看成“工具入门”。它训练的是最终项目最底层的可信证据：别人能不能找到你的代码、恢复你的环境、确认你的数据入口，并理解你的图表为什么能支持某个有限论断。

\section*{AI 助手在 Part I 之后应该怎样使用}

到了这里，AI 助手可以帮助你写命令、解释报错、整理脚本、检查 README、生成图表说明和提醒复现清单。但有三件事不能外包：

\begin{itemize}
    \item 你必须知道当前命令会读写哪些文件；
    \item 你必须确认数据字段、单位和来源是否可信；
    \item 你必须能把图表和结果限制讲给别人听。
\end{itemize}

换句话说，AI 助手可以提高工作速度，但不能替你承担科研记录的责任。越是使用 AI，越要保留清楚的运行证据、版本证据和判断边界。

\section*{进入 Part II 前}

Part II 会开始处理具体天文和物理数据：误差、星表、图像、光谱、时间序列和实验模型。那时你会发现，很多“数据分析问题”其实会先退回 Part I 的地基问题：路径是否稳定、数据是否有契约、代码是否能重跑、图是否能支撑解释、结果是否能对应到某次提交。

如果这些问题已经能回答，Part II 就不会只是学习更多文件格式，而会自然变成下一步：把一个可复现项目包里的数据，读成带物理含义的科学证据。

\section*{一个可以立刻执行的 Part I 综合任务}

如果读者想确认自己是否真的准备好进入 Part II，可以把五章合并成一个两小时以内的小任务：

\begin{enumerate}
    \item 创建一个最小项目目录 \texttt{mini\_research\_package/}，至少包含数据、脚本、notebook、结果、图表和日志六类子目录；
    \item 用命令行生成目录快照和输入文件清单；
    \item 初始化 Git 仓库，提交 README、脚本和数据说明；
    \item 写一个 Python 脚本，读取小型 CSV，检查字段和缺失值，输出清洗后摘要；
    \item 生成一张带单位、图注和 claim--evidence--limit 的图；
    \item 写一份 evidence record，包含命令、环境、commit hash、输入、输出、数据限制和图表解释；
    \item 打一个阶段性 tag，说明这份包可以支持什么结论、不能支持什么结论。
\end{enumerate}

这个任务不需要复杂数据，也不需要高级模型。它检查的是最基础但最容易被忽略的能力：一个结果能否被找到、复跑、追溯、检查和解释。能完成它，才说明 Part I 不只是读过，而是真的成为了后续学习的地基。
